\section{Existence and Uniqueness of Weak Solutions}

In this section, we establish the well-posedness of the stochastic advection--diffusion problem introduced in the previous section.

\subsection{Continuity of the Bilinear Form}

Recall that the bilinear form is defined by

\[
a(u,v)
=
\kappa(\nabla u,\nabla v)
+
(\mathbf b\cdot\nabla u,v),
\]

for all $u,v\in V$.

\begin{lemma}
The bilinear form $a(\cdot,\cdot)$ is continuous on $V\times V$. More precisely, there exists a constant $C>0$ such that

\[
|a(u,v)|
\le
C
\|u\|_V
\|v\|_V
\]

for all $u,v\in V$.
\end{lemma}

\begin{proof}
Using the Cauchy--Schwarz inequality,

\[
|\kappa(\nabla u,\nabla v)|
\le
\kappa
\|\nabla u\|
\|\nabla v\|.
\]

Moreover,

\[
|(\mathbf b\cdot\nabla u,v)|
\le
\|\mathbf b\|_{L^\infty}
\|\nabla u\|
\|v\|.
\]

Applying the Poincaré inequality,

\[
\|v\|
\le
C_P
\|\nabla v\|,
\]

we obtain

\[
|a(u,v)|
\le
C
\|u\|_V
\|v\|_V.
\]

Hence the result follows.
\end{proof}

\subsection{Coercivity}

\begin{lemma}
Assume that

\[
\nabla\cdot\mathbf b=0.
\]

Then the bilinear form $a(\cdot,\cdot)$ is coercive on $V$:

\[
a(v,v)
\ge
\kappa
\|\nabla v\|^2
\]

for all $v\in V$.
\end{lemma}

\begin{proof}

We have

\[
a(v,v)
=
\kappa\|\nabla v\|^2
+
(\mathbf b\cdot\nabla v,v).
\]

Since $\nabla\cdot\mathbf b=0$ and $v=0$ on $\partial\Omega$,

\[
(\mathbf b\cdot\nabla v,v)=0.
\]

Therefore,

\[
a(v,v)
=
\kappa
\|\nabla v\|^2.
\]

Since $\kappa>0$, the result follows.
\end{proof}

\subsection{Existence of Weak Solutions}

We are now in a position to establish the existence of a weak solution.

\begin{theorem}
Assume that

\[
u_0\in L^2(\Omega),
\]

\[
f\in L^2(0,T;L^2(\Omega)),
\]

and

\[
\mathbf b\in [L^\infty(\Omega)]^2.
\]

Then there exists a weak solution

\[
u
\in
L^2\bigl(
\Omega_P;
L^2(0,T;V)
\bigr)
\]

to the stochastic advection--diffusion problem.
\end{theorem}

\begin{proof}[Sketch of proof]

The proof is based on the Galerkin approximation method.

Let

\[
V_m
=
\operatorname{span}
\{\varphi_1,\ldots,\varphi_m\}.
\]

be a finite-dimensional subspace of $V$.

We seek an approximation

\[
u_m(t)
=
\sum_{i=1}^{m}
c_i(t)\varphi_i.
\]

The coefficients satisfy a finite-dimensional stochastic system.

Using the a priori estimate established in Section 3, one obtains uniform bounds with respect to $m$.

Compactness arguments and weak convergence allow passage to the limit.

Hence a weak solution exists.
\end{proof}

\subsection{Uniqueness}

\begin{theorem}
The weak solution of the stochastic advection--diffusion problem is unique.
\end{theorem}

\begin{proof}

Let $u_1$ and $u_2$ be two weak solutions.

Define

\[
w=u_1-u_2.
\]

Subtracting the two equations yields

\[
dw
=
Aw\,dt.
\]

Testing with $w$ gives

\[
\frac12
\frac{d}{dt}
\|w\|^2
+
a(w,w)
=
0.
\]

Using coercivity,

\[
\frac12
\frac{d}{dt}
\|w\|^2
+
\kappa
\|\nabla w\|^2
=
0.
\]

Since

\[
w(0)=0,
\]

we obtain

\[
\|w(t)\|^2=0
\]

for all $t\in[0,T]$.

Therefore

\[
u_1=u_2.
\]

Hence the solution is unique.
\end{proof}